\documentclass[a4paper]{article}
\usepackage[utf8]{inputenc}
\usepackage[T1]{fontenc}
\usepackage{lmodern}
\usepackage[ngerman]{babel}
\usepackage{amsmath,amssymb}
\usepackage{graphicx}
\usepackage{booktabs}
\usepackage{siunitx}
\usepackage{listings}
\usepackage{makeidx}
\usepackage[numbers]{natbib}
\usepackage{hyperref}
\makeindex
\title{Ein Testpapier für den Amiga}
\author{Jörg Müller}
\begin{document}
\maketitle
\begin{abstract}Größere Umlaute: äöüß ÄÖÜ. Silbentrennung: Donaudampfschifffahrtsgesellschaftskapitän.\end{abstract}
\tableofcontents
\section{Einleitung}\label{sec:intro}
Siehe \citet{knuth84} und \cite{lamport86,mueller20}. \index{Knuth}\index{Amiga!Vampire}
Abschnitt~\ref{sec:math} auf Seite~\pageref{sec:math}.
\begin{figure}[ht]\centering\includegraphics[width=4cm]{img}\caption{Ein Bild}\end{figure}
\section{Mathematik}\label{sec:math}
\begin{align}
  \nabla\cdot\mathbf{E} &= \frac{\rho}{\varepsilon_0} \\
  \int_{-\infty}^{\infty} e^{-x^2}\,dx &= \sqrt{\pi}
\end{align}
Die Geschwindigkeit beträgt \SI{3.0e8}{\metre\per\second}.\index{Geschwindigkeit}
\begin{table}[ht]\centering
\begin{tabular}{lrr}\toprule Modell & CPU & RAM \\\midrule A1200 & 68020 & \SI{2}{\mega\byte} \\ V4SA & 68080 & \SI{512}{\mega\byte} \\\bottomrule\end{tabular}
\caption{Rechner}\end{table}
\begin{lstlisting}[language=C,basicstyle=\ttfamily\small]
int main(void) { printf("Hello\n"); return 0; }
\end{lstlisting}
\bibliographystyle{plainnat}
\bibliography{refs}
\printindex
\end{document}
